\chapter{What is life?}

\section{Why the philosophy chapter is here at all}

A book about artificial life that ducks the question \emph{what is life}
will get away with it for a while and then trip. Sooner or later you will
write down a viability filter or an observer and somebody will ask why
\emph{that} one and not another, and you will have to either appeal to a
prior commitment about what matters or pretend the choice was arbitrary.
This chapter equips you with the prior commitments. It surveys the frames
that have shaped the field --- Aristotle, the mechanists, Schr\"odinger,
the autopoieticists, the active-inference theorists, and the
information-theoretic individualists --- and shows where each of them lives
inside the five-slot lens introduced in Chapter~4.

The chapter is not a defence of any one frame. Several of the frames are
mutually incompatible at the level of metaphysics, and the book will not
ask you to pick one. The point is to make the choice \emph{visible}: when
you read a paper that talks about ``the agent'' or ``the organism,'' you
should be able to spot which frame the author is borrowing from, and you
should be able to say, in the vocabulary of the lens, what slot the frame
is doing work in.

If this chapter feels like a history of ideas in twenty pages, that is
deliberate. The deep technical work begins in Chapter~4. What we are doing
here is laying the conceptual ground for it.

\section{Aristotle: life as form}

Aristotle's account of life is older than the word \emph{biology} and still
hard to dismiss. He observed that living things --- plants, animals,
people --- differ from rocks not by being made of different stuff but by
being organised in a particular way. A corpse and a person are made of the
same materials in the same arrangement, momentarily, and yet only one of
them is alive. What distinguishes them, Aristotle argued, is the
\emph{form} that organises the matter: the soul (\textit{psych\^e}), which
in his usage is not a ghost but a principle of self-organising activity.

The Aristotelian living thing has four causes --- material, efficient,
formal, and final. The material cause is what it is made of. The efficient
cause is what made it. The formal cause is what it is configured as. The
final cause is what it is for. Plants have nutritive souls (they grow and
reproduce). Animals add sensitive souls (they perceive and move). Humans
add rational souls. The hierarchy was tidy. None of it survived the
scientific revolution intact, but two pieces did: the idea that life is a
matter of \emph{organisation} rather than substance, and the idea that
some living systems are organised around a \emph{purpose} that the parts
serve.

For our purposes, Aristotle's lasting contribution is the move he makes in
the first place: separating what something is made of from what it does.
In the language of the five-slot lens, this maps almost directly onto the
separation between substrate (what the system is made of) and the
operators that act on it (variation, viability, topology, observer). A
Lenia blob is made of a real-valued field on a lattice. The fact that we
can talk about it as a \emph{creature} --- something that grows, persists,
divides --- is not a fact about the lattice. It is a fact about the form
imposed on the lattice by the update rule and the persistent pattern that
the rule admits.

What Aristotle does not give us is a way of telling whether a particular
arrangement of matter counts as alive. That problem will be with us for
the rest of the book.

\section{The mechanists: life as a kind of machine}

The mechanist tradition, beginning with Descartes in the seventeenth
century and culminating in nineteenth-century physiology, made the
opposite move. Where Aristotle distinguished living things from non-living
by an extra organising principle, the mechanists asserted that there was
no extra principle. Living things were machines. Complicated ones, with
parts that interacted in subtle ways, but machines all the same. Descartes
famously argued that animals were automata; he was less brave about
extending the claim to humans, but his contemporaries were less hesitant.

The mechanist programme paid off scientifically. Harvey's account of
circulation, the discovery of the cell, the chemistry of metabolism, the
identification of DNA --- each of these found a mechanism in a place where
earlier writers had been content with a principle. The success of the
programme is the reason the question \emph{what is life} can still feel
slightly embarrassing. We know a great deal about how living things work.
What more is there to ask?

The thing the mechanists left on the table, and that the artificial-life
community returned to, is the question of which organisations of mechanism
are alive in the first place. Descartes did not propose a criterion for
distinguishing a clockwork from a clockwork-that-is-alive. The criterion
he gestured at --- complexity --- did not survive contact with twentieth
century cellular biology, which discovered that the cell is staggeringly
more complex than the most elaborate machine of the mechanist period and
yet, from the mechanist's perspective, just another machine.

The lens we will build does not need to settle the metaphysics. It will
treat the substrate as a machine in the mechanist sense (a measurable
space with a measurable update rule), and ask separately, through the
viability and observer slots, what kinds of organisations on the machine
we are willing to call alive.

\section{Schr\"odinger: life as what resists dissipation}

In 1944 Erwin Schr\"odinger, then in exile in Dublin, gave a series of
lectures at Trinity College titled \textit{What is Life?}. He was not a
biologist. He came to the question from physics, and his physics-flavoured
answer turned out to be unusually productive.

Schr\"odinger's puzzle was the second law of thermodynamics. Closed systems
trend toward equilibrium; living things appear to do the opposite. A
mouse is a far-from-equilibrium structure that maintains itself for years.
How is this consistent with the second law? Schr\"odinger's answer was
that organisms are not closed --- they pump entropy out into the
environment. ``[Life] feeds upon negative entropy,'' he wrote, by which he
meant that organisms maintain their internal order by exporting disorder
to the outside.

The phrase is rough but the idea is exact. A living thing is a structure
through which energy flows in a way that maintains a far-from-equilibrium
configuration. Stop the flow, and the configuration relaxes. The metaphor
of life-as-flame catches this: a flame is a self-sustaining configuration
of combustion, complete with apparent boundaries (the silhouette of the
flame) and apparent agency (the way it ``reaches'' for fuel). It is not
alive, but the structural similarity is informative.

Schr\"odinger also speculated about a ``code-script'' in the chromosome ---
years before the discovery of DNA --- that would carry the information
needed to specify the organism's form. He thought of it as an
``aperiodic crystal'': structured enough to encode something, not periodic
enough to repeat. He was right in the spirit if not the letter; molecular
biology would make the picture concrete a decade later.

Schr\"odinger's main contribution to our vocabulary is the idea that life
is a \emph{persistence} phenomenon. A living thing is something that
keeps being itself, against a thermodynamic drift toward not being
itself. In the five-slot lens, this maps onto the viability slot --- the
filter that decides which configurations persist and which fade --- and
we will return to it in Chapter~6.

\section{Autopoiesis: life as self-production}

Schr\"odinger told us that living things resist dissipation. He did not
tell us \emph{how}. The Chilean biologists Humberto Maturana and Francisco
Varela, working in the 1970s, proposed an answer: a living system is one
whose components produce, and are produced by, each other. They called
this structure \emph{autopoiesis} --- self-production.

An autopoietic system has a network of processes (chemical, in the
biological case) such that:
\begin{itemize}[leftmargin=1.5em,itemsep=2pt]
  \item each process in the network is produced by other processes in the
    network;
  \item the network as a whole produces the boundary that separates it
    from its environment;
  \item that boundary is what makes the network identifiable as a unit.
\end{itemize}

A cell is the canonical example. The cell membrane is produced by the
chemistry of the cell; the chemistry of the cell is what the membrane
sustains; and the membrane is what makes the cell a discrete thing rather
than a smear of organic chemistry. Take any one of these away and the
cell is no longer recognisable as a cell.

Autopoiesis is a closure condition. The set of processes is closed under
the relation of mutual production --- nothing inside the set requires
anything outside the set for its production. This formal feature will
matter in Chapter~6: it is one of the four ways of being precise about
\emph{which configurations persist}. The Maturana--Varela framework was
articulated in biological language, but Randall Beer has shown in a long
series of papers that gliders in Conway's Game of Life satisfy a
non-trivial autopoietic closure condition --- they are, in a precise
sense, autopoietic units of the GoL ``reaction network.'' That is one of
the reasons the Game of Life keeps coming up in artificial-life
arguments.

The autopoietic frame is not the only formalisation of life-as-closure.
The RAF set theory of Kauffman, Hordijk, and Steel gives another: a set
of reactions in a chemical network is reflexively autocatalytic and
food-generated (a RAF set) if every reaction in it is catalysed by
something in it (or in a food set), and every reactant is producible
from the food set inside the RAF. RAFs and autopoiesis are not the same
notion, but they share the shape: the system is closed under a relevant
relation. Whether they are \emph{secretly} the same notion, in the sense
that all of the viability-as-closure proposals collapse to a common
closure operator, is an open conjecture (C1 in the paper; see
\texttt{RESULTS.md}). The book takes the conjecture seriously and does
not assume it.

\section{Active inference: life as resisting surprise}

Karl Friston's free energy principle (FEP) is the most ambitious modern
candidate for a unifying theory of life and cognition. Its claim, in
slogan form: any system that persists over time must act so as to minimise
the long-run surprise of its sensory inputs. Equivalently, it must behave
as if it were a Bayesian inference engine, continually updating its
beliefs about the world and acting to keep its actual sensory states
consistent with what its model predicts.

The technical apparatus is dense. We will summarise rather than derive.
Friston starts from a Markov blanket --- a statistical decomposition of a
system into internal states, sensory states, active states, and external
states, such that internal and external states are conditionally
independent given the blanket (sensory + active) --- and shows that, under
certain conditions, the internal states can be interpreted as
parameterising a generative model whose free energy they minimise. A
thermostat does this trivially. A bacterium does it less trivially. A
human, in Friston's telling, does it most non-trivially of all.

The FEP has been received with both enthusiasm and skepticism. The
enthusiasm is for its breadth: it offers a single mathematical principle
that ostensibly accounts for action, perception, learning, and
homeostasis. The skepticism is for the same reason. Aguilera, Biehl, and
collaborators have argued that the Markov blanket decomposition is harder
to find in real systems than the slogans suggest, and that the bridge
from sparse coupling to genuine blanket structure requires conditions
that are usually not checked. Friston's reply sharpens but does not
resolve the dispute. The book takes the FEP as a \emph{descriptive lens}
in the same spirit as the rest of the framework: when a system admits a
Markov blanket decomposition that is preserved under its dynamics, you
may parameterise its internal dynamics as free-energy minimisation. When
it does not, you may not. The question of when it does is open.

In the lens, the FEP sits across two slots. The Markov blanket condition
is a viability formalism (one of the four in Chapter~6). The
free-energy-minimising interpretation of internal-state dynamics is
substrate-side: it parameterises $\Phi$ for systems where it applies. The
FEP does not occupy its own slot. It occupies particular configurations
of the slots already in the lens.

\section{Information-theoretic individuality: life as predictive closure}

The newest entry in this chapter is information-theoretic individuality.
Krakauer, Bertschinger, Olbrich, Flack, and collaborators have proposed
that an individual --- the part of a system that we call ``a thing'' ---
is characterised by a kind of \emph{informational closure}. Roughly: the
future behaviour of an individual can be predicted from its own past
better than from the past of the world outside it.

The formalisation runs through mutual information. Let $S$ be a candidate
``inside'' (a set of substrate sites). Let $\bar S$ be the complement
(the outside). Roughly, $S$ counts as an individual if the information
that the past of $S$ carries about its own future is large compared to
the additional information the outside carries, conditional on the
inside. (If the phrase \emph{mutual information} is new, the math primer
at the back of the book is the gentle introduction; don't worry about
the precise formula yet.)

There are several variants. The book uses one in Chapter~4's discussion
of entities. The variants share a structural commitment: \emph{the inside
of a thing is informationally distinguishable from its outside}. This is
not the same as the autopoietic closure --- a system can be
informationally closed without being self-producing, and vice versa ---
but the two are related, and how related is something the framework
leaves open.

For the lens, this frame supplies the entity definition. An entity, in
the paper's sense, is a supporting set $S$ (the sites that hold it) plus
a scale projection $\pi$ (the description we are committing to). The
informational closure condition is what makes $S$ count as an entity in
the first place. We will unpack this in Chapter~4 once the lens is on the
table.

\section{Where the frames land in the lens}

Five frames, five chapters of intellectual history, and we have been
careful not to take a side. The reason will be visible if you tabulate
where each frame does its work:

\begin{center}
\begin{tabular}{@{}p{0.30\textwidth}p{0.61\textwidth}@{}}
\hline
\textbf{Frame} & \textbf{What slot(s) it speaks to} \\
\hline
Aristotle (form over matter) &
  Separation of substrate from operators; entity scale projection. \\
Mechanism (life is machinery) &
  Substrate dynamics. \\
Schr\"odinger (life as persistence) &
  Viability filter. \\
Autopoiesis (life as self-production) &
  Viability (autopoietic-closure variant). \\
RAF (catalytic closure) &
  Viability (RAF variant). \\
Active inference / FEP &
  Viability (Markov-blanket variant); substrate dynamics
  (free-energy parameterisation). \\
Info-theoretic individuality &
  Entity definition (supporting set $S$ and projection $\pi$). \\
\hline
\end{tabular}
\end{center}

Notice that the philosophical frames do not collide. They sit in
different parts of the same structure. The mechanist's machine is the
substrate; the autopoieticist's closure is one of several viability
filters; the active-inference theorist's free-energy minimisation
parameterises the substrate dynamics; the information-theoretic
individualist supplies the entity definition. The lens does not arbitrate
between them, because they are not really fighting --- they are
describing different parts of the same animal.

That is the trick the rest of the book runs. The slots are a place to
put each candidate answer to ``what is life?'' so that we can see them
side by side, choose among them when we have to, and identify what each
one is committing us to.

\section{Two warnings before we leave the philosophy chapter}

\textit{Warning 1: the slots are not magic.} The fact that a frame ``fits
in the lens'' does not mean the frame is correct. The free energy
principle, for example, fits cleanly into the viability slot --- and yet
the empirical question of whether real organisms admit the Markov
blanket structure the principle requires is unsettled. The lens lets you
write down a frame's commitment; it does not certify the commitment.

\textit{Warning 2: artificial life does not need a settled
metaphysics.} You can build Lenia, Boids, Tierra, or POET without
choosing between Aristotle and Friston. The lens is intentionally
descriptive, and you will spend most of the book in this mode: reading
systems, comparing them, writing them down. The deep questions of what
life \emph{is} will remain in the background. The slots make it possible
to do useful work while those questions are unresolved.

\section*{To think about}

\begin{enumerate}
  \item Pick a system you know well (your operating system's process
    scheduler, an ant colony, a bacterial culture). Which of the frames
    in this chapter would you reach for to describe it? Which slot does
    that frame occupy?
  \item Schr\"odinger's ``life feeds on negative entropy'' is sometimes
    presented as if it solved the puzzle of life. Why is it not a
    solution, but a constraint?
  \item Autopoiesis and RAF are both closure conditions on a network of
    processes. Find one biological example where the autopoietic closure
    is non-trivial but the RAF closure is trivial, and one where the
    reverse holds.
  \item The free energy principle is sometimes stated as a tautology
    (``anything that persists must do this'') and sometimes as an
    empirical claim (``brains do this''). Which version does the lens
    accommodate?
  \item Pick a Lenia creature. Write down, in plain prose, what its
    autopoietic boundary would be. Where does the creature stop and the
    field begin?
\end{enumerate}
